\documentclass[11pt]{article}

% font
\usepackage{tgschola}

\usepackage{graphicx}
\usepackage[T1]{fontenc}
\usepackage[utf8]{inputenc}
\usepackage[english]{babel}
\usepackage{multicol}
\usepackage[ddmmyyyy]{datetime}
\usepackage[a4paper, margin=3cm]{geometry}

% sections on new page
\usepackage{titlesec}
\newcommand{\sectionbreak}{\clearpage}
\setcounter{secnumdepth}{0}

% index and clickable toc
\usepackage{hyperref}
\hypersetup{
	colorlinks,
	citecolor=black,
	filecolor=black,
	linkcolor=black,
	urlcolor=black
}

% table of content name
\addto\captionsenglish{% Replace "english" with the language you use
	\renewcommand{\contentsname}%
	{\hfill \Huge Table of Content \hfill}%
}

\begin{document}

	% Title page
	\begin{titlepage}
		\begin{center}

			\vspace*{0.5cm}

			% title
			\Huge
			\textbf{Brain Tumor Detection}

			\vspace{0.8cm}

			% logo
			\includegraphics[width=0.4\linewidth]{imgs/cu_logo.png}

			\vspace{1cm}

			\normalsize
			Department of Computer Science\\
			Gurudas College\\
			Calcutta University\\
			\today

			% subtitle
			\vspace{0.5cm}
			\textit{Detection of tumorous cells using machine learning models}

			\vspace{1cm}

			\normalsize
			Submitted in partial fulfillment for the requirements for the
			degree Bachelor of Science (Honors) in Computer Science.

			Academic Year: \textbf{2018-2021}

			\vfill
			% author
			\textbf{Authors}:

			$\bullet$ Shoptorshi \space\space
			$\bullet$ Brahmajit \space\space
			$\bullet$ Rajarshi \space\space
			$\bullet$ Bhargav

			% supervisor
			\vspace{1cm}
			\rule{5cm}{0.4pt}

			\textbf{Supervisor}: Kajari Bhattacharjeee

		\end{center}
	\end{titlepage}

	% certificate page
	\section[certificate]{\hfill \Huge Certificate \hfill}
	%\thispagestyle{empty}

	% roman page numbering
	\pagenumbering{roman}

	\begin{center}
		\vspace{1cm}

		% logo
		\includegraphics[width=0.3\linewidth]{imgs/cu_logo.png}

		\vspace{1cm}

		\Large
		Department of Computer Science\\
		Gurudas College\\
		Calcutta University\\
	\end{center}

	\vspace{1.8cm}
	This is to certify that the project entitled "Brain tumor detection using
	Machine Learning models" is a bona fide work of \textbf{Shoptorshi},
	\textbf{Brahmajit}, \textbf{Rajarshi} and \textbf{Bhargav} submitted to
	Gurudas College, University of Calcutta; in partial fulfillment of the
	requirement for the award of the degree "Bachelor of Science (Honors)" in
	Computer Science.
	\vfill

	\begin{multicols}{2}

		Supervisor:

		\rule{5cm}{0.4pt}

		\textbf{Kajari Bhattacharjeee}

		\vspace{1cm}

		Department Head:

		\rule{5cm}{0.4pt}

		\textbf{Srijeeta Charkraborty}


		\columnbreak

		Principal:

		\rule{5cm}{0.4pt}

		\textbf{Dr. Mausumi Chatterjee}

	\end{multicols}

	\section[acknowledgement]{\hfill \Huge Acknowledgement \hfill}
	\label{sec:Acknowledgement}


	I would like to express my special thanks of gratitude to our supervisor
	\textbf{Kajari Bhattacharjeee}, teachers of our department \textbf{Sonali
	Gupta} and department head \textbf{Srijeeta Charkraborty}; as well as our
	principal ma'am \textbf{Dr. Mausumi Chatterjee} who gave us the golden
	opportunity to do this wonderful project on \textit{Brain tumor detection},
	which has also helped us in doing a lot of research and we came to know
	about many new things.

	We are really thankful to them.

	Secondly we would like to thank the incredible authors of the research
	papers in our references/citations for providing us with the information.

	\tableofcontents

	% abstract
	\section[abstract]{}

	% reverting back to normal page numbering
	\pagenumbering{arabic}

	\thispagestyle{plain}
	\begin{center}
		\Large
		\textbf{Brain Tumor Detection}

		\vspace{0.4cm}
		\large
		Using various machine learning models to detect brain tumor.

		\vspace{0.5cm}
		\textbf{Abstract}
	\end{center}

	Tumors are cancerous or non-cancerous mass or growth of abnormal cells in
	brain. Tumors can start in brain, or cancer elsewhere in the body can spread
	to brain. There are many way to control the occurrence of these abnormal
	cells. A tumor can be denoted as a malformed mass of tissues wherein the
	cells multiply abruptly and ceaselessly, that is there is no control over
	the growth of the cells.

	The process of Image segmentation is adopted for extracting abnormal tumor
	region within the brain. In the MRI (magnetic resonance image), segmentation
	of brain tissue holds very significant in order to identify the presence of
	outlines concerning the brain tumor. There is abundance of hidden
	information in stored in the Health care sector. With appropriate use of
	accurate data mining classification techniques, early prediction of any
	disease can be effectively performed.

	The project examines list of risk factors that are being traced out in brain
	tumor surveillance systems. Also the method proposed assures to be highly
	efficient and precise for brain tumor detection, classification and
	segmentation. To achieve this precise automatic or semi-automatic methods
	are needed. The project proposes an automatic segmentation method that
	relies upon \textit{ CNN (Convolution Neural Networks) }, \textit{ VGG 16 }
	and \textit{ Resent 50 }, determining small 7 x 7 kernels. By incorporating
	this single technique, segmentation and classification is accomplished. CNN
	(a ML technique) from NN (Neural Networks)wherein it has layer based for
	results classification.

	Various levels involved in the proposed mechanisms are:

	\begin{enumerate}
		\item \textbf{ Data collection }
		\item \textbf{ Pre-processing }
		\item \textbf{ Average filtering }
		\item \textbf{ segmentation }
		\item \textbf{ feature extraction }
		\item \textbf{CNN (or any other model) via classification and identification. By
			utilizing the DM (data mining) techniques, significant relations and
			patterns from the data can be extracted. The techniques of ML
			(machine learning) and Data mining are being effectively employed
			for brain tumor detection and prevention at an early stage.}
	\end{enumerate}

	\section[introduction]{\hfill \Huge Introduction \hfill}
	\subsection[domain description]{Domain Description}

	\begin{enumerate}
		\item \textbf{Neurological Examination}: It is a series of test to
			measures the function of the patients nervous system and also
			his/her physical and mental alertness.
		\item \textbf{Machine Learning}: Machine learning approaches address
			these problems by mainly using hand-crafted features (or pre-defined
			features). As an initial step in this kind of segmentation, the key
			information is extracted from the input image using some feature
			extraction algorithm, and then a discriminative model is trained to
			recognize the tumor from normal tissues. The designed machine
			learning techniques generally employ hand-crafted features with
			various classifiers, such as random forest, support vector machine
			(SVM), fuzzy clustering. The designed methods and features
			extraction algorithms have to extract features, edge-related
			details, and other necessary information—which is time-consuming.
			Moreover, when boundaries between healthy tissues and tumors are
			fuzzy/vague, these methods demonstrate poorer performances.
		\item \textbf{Brain Scan}: Brain scan is a picture of the internal
			structure of the brain. A specialized machine takes a scan in the
			same way as a digital camera takes a photograph. Using computer
			technology, a scan compiles an image of the brain by photographing
			it from various angles. Some types of scan uses contrast agent (or
			contrast dye), which helps the doctor to see the difference between
			normal and abnormal brain tissues.

			MRI (Magnetic Resonance Imaging): It is a scanning device that uses
			magnetic field and computer to capture images of the brain on films.
			It does not use x-rays. It provides pictures from various planes,
			which permits doctor to create a three-dimensional image of the
			tumor. The MRI detects signals emitted from normal and abnormal
			tissues, providing clear images of almost all tumors.
	\end{enumerate}

	\subsection[motivation]{Motivation}
	The motivation is to develop a software with better segmentation capability
	for use in medical imaging to detect diseases like brain tumor. Image
	segmentation has been identified as the key problem of medical image
	analysis and remains a popular and challenging area of research. Image
	segmentation is increasingly used in many clinical and research applications
	to analyze medical imaging datasets; which motivated us to present a
	snapshot of dynamically changing field of medical image segmentation.

	CT (Computed Tomography), MRI (Magnetic Resonance Imaging), PET (Positron
	Emission Tomography) etc. generates a large amount of image information.
	With the improved technology, not only does the size and resolution of the
	images grow but also the number of dimensions increases. In the future, we
	would like to have algorithms which can automatically detect diseases,
	lesions and tumors, and highlight their locations in the large pile of
	images.

	The motivation of this work is to increase patient safety by providing
	better and more precise data for medical decision.

\end{document}
